\documentclass[11pt,a4paper]{article}
\usepackage[utf8]{inputenc}
\usepackage[T1]{fontenc}
\usepackage{lmodern}
\usepackage[margin=1in]{geometry}
\usepackage{amsmath,amssymb,amsthm}
\usepackage{booktabs}
\usepackage{array}
\usepackage{microtype}
\usepackage[hidelinks]{hyperref}

\newtheorem{theorem}{Theorem}
\newtheorem{corollary}[theorem]{Corollary}
\theoremstyle{remark}
\newtheorem{remark}[theorem]{Remark}

\newcommand{\tier}[1]{{\small\textsf{[#1]}}}
\newcommand{\src}[1]{{\footnotesize\texttt{#1}}}
\newcommand{\Ad}{A_\delta}
\newcommand{\Sd}{S_\delta}

\title{\bfseries Parity Blindness of Finite-Window Weil Detectors at a Real Zero:
the Even Sector Cannot See a Siegel Configuration and the Odd Sector Is Maximally
Sensitive to It}

\author{Elias De Jes\'us\\[2pt]
{\small Independent Researcher}\\[1pt]
{\small ORCID \href{https://orcid.org/0009-0007-0190-9143}{0009-0007-0190-9143}}\\[1pt]
{\small \texttt{dejesuselias10@gmail.com}}}

\date{October 2026 --- Technical Note\\[2pt]
{\small\texttt{doi:10.5281/zenodo.23623473}}}

\begin{document}
\maketitle

\begin{abstract}
\noindent
In the finite-window Weil form, a real zero pair $\{\tfrac12+\delta,\tfrac12-\delta\}$
contributes $2F(i\delta)F(-i\delta)$, and because the autocorrelation factorises this is
$+2\Ad(f)^2$ for even test functions and $-2\Sd(f)^2$ for odd ones, with
$\Ad=\int f\cosh(\delta u)\,du$ and $\Sd=\int f\sinh(\delta u)\,du$. The perturbation is
therefore rank one, positive semidefinite on the even sector and negative semidefinite on
the odd sector, for every window length and in every basis. Two consequences follow. Even-sector
finite-window positivity is compatible with a real zero at every $\delta\in(0,\tfrac12)$,
so no even-sector positivity computation --- at any $L$, in any basis, to any precision ---
carries information about the Siegel configuration; the window is not merely blind there,
it is anti-detecting, and the statement is exact in $\delta$ rather than a vanishing
leading coefficient. Conversely the odd-sector sensitivity at $\gamma=0$ is the
$\gamma$-uniform maximum $2\cdot(2L^3/3)$ of the identity
$K_{\mathrm{even}}(L,\gamma)+K_{\mathrm{odd}}(L,\gamma)=2L^3/3$, so $\gamma=0$ is the
unique ordinate at which the two sectors' roles separate completely. The sign structure is
verified in $150$ cells to $8.49\times10^{-61}$, reproduces the repository's independently
derived odd-sector pole term entrywise with no normalisation adjustment, and holds exactly
over $\mathbb{Q}$ in the function-field control. This resolves non-claim~9 of the companion
note~\cite{DHnote}.
\end{abstract}

\section*{0.\quad Non-claims}
\addcontentsline{toc}{section}{Non-claims}

Stated first.

\begin{enumerate}\itemsep3pt
\item \textbf{Nothing about RH or GRH}, and \textbf{nothing about whether Siegel zeros
exist.} The theorem below describes what a window form does when a real zero is placed in
the zero multiset by hand. It is a statement about a detector, not about the zeros of any
$L$-function.

\item \textbf{The quasi-Riemann-hypothesis and Siegel-exclusion preprints of October 2026
are not cited as established and are not used anywhere.} No result here depends on them,
and none is asserted about them.

\item \textbf{Every planted detection is \tier{T2-N}.} A planted real zero is a
perturbation of a computed matrix, not an arithmetic fact, and
\textbf{non-detections certify nothing}.

\item \textbf{``The odd sector detects by $L=1.4$'' is \tier{T2-N} twice over.} It is a
numerical statement about a planted configuration, and the odd-sector \emph{baseline} it is
measured against has no certified lower bound: no Zhu-style certificate exists for any
Dirichlet window. Converting the measurement into an exclusion of a real zero would require
exactly such a bound, and this note supplies none.

\item \textbf{The parent extension's Dirichlet-character results beyond its \S5 are not
used here.} Saturation ratios, prime-side coincidences, conductor-versus-height
observations and the Phase-4 mechanism hypothesis play no part in anything below.

\item \textbf{Corollary~\ref{cor:even} bounds what a method can detect, not what exists.}
Section~\ref{sec:lit} is a scope statement about even-sector techniques. It is not a
criticism of any result, and no cited work claims otherwise.

\item \textbf{The function-field section is a control, not a transfer.} In the function
field the analogue of RH is Weil's theorem, so there is no configuration there to exclude.
\end{enumerate}

\section*{Epistemic tiers}
\addcontentsline{toc}{section}{Epistemic tiers}

\begin{center}
\begin{tabular}{@{}lll@{}}
\toprule
Tier & Meaning & Discharged by \\
\midrule
T1 & exact or proved & a proof, or exact arithmetic over $\mathbb{Z}$ or $\mathbb{Q}$ \\
T2 & numerical & stated working precision $+$ a convergence or noise-floor check \\
T2-C & certified arithmetic & an interval/ball enclosure with a reproducible certificate \\
T2-N & high-precision numerical & stable under precision increase, not an enclosure \\
T3 & interpretation & argument from T1/T2 evidence, labelled as such \\
\bottomrule
\end{tabular}
\end{center}

\noindent Every number printed in this note is re-derived from the extension's data files
by \src{note/verify\_note\_numbers.py}, which also fails if the note prints more
significant digits than its source carries.

\section{Setup}\label{sec:setup}

Let $g$ be real and supported in $[-L,L]$. Following the companion note~\cite[\S2]{DHnote},
put
\[
G(z)=\int g(u)e^{zu}\,du,\qquad h(v)=\int g(u+v)g(u)\,du,\qquad H(z)=G(z)G(-z),
\]
so $h$ is supported in $[-2L,2L]$ and $H(it)=|G(it)|^2$ for real $g$. The spectral side of
the explicit formula runs over the zeros of the completed function, and all reported values
are normalised as $Q(g)/\|g\|_2^2$. For $\zeta$ the window form is this repository's
assembled $Q_\zeta$; for a real primitive $\chi$ modulo $q$ it is
\[
M_\chi=\mathrm{Arch}_a-\mathrm{LogPi}+\mathrm{LogQ}-\mathrm{Prime}_\chi ,
\]
with no pole term, $\chi$ being non-principal. Both parity sectors are used throughout:
the even bases $\cos((2k+1)\pi x/2L)$ and $\cos(k\pi x/L)$, and the odd basis
$\sin(k\pi x/L)$, all normalised.

For a non-real off-line zero the complete functional-equation and conjugation orbit is
$\mathcal{O}=\{\beta\pm i\gamma,\;1-\beta\pm i\gamma\}$, and its contribution
is~\cite[eq.~(4)]{DHnote}
\begin{equation}\label{eq:orbit}
Q_{\mathcal O}(g)=4\,\mathrm{Re}\!\left(A\overline{B}\right)=4\left(|P|^2-|S|^2\right),
\qquad A=G(\delta+i\gamma),\quad B=G(-\delta+i\gamma),
\end{equation}
with $P=\int g\cosh(\delta u)e^{i\gamma u}du$ and $S=\int g\sinh(\delta u)e^{i\gamma u}du$,
$\delta=\beta-\tfrac12$: a signed rank-two form, positive on the $\cosh$ channel and
negative on the $\sinh$ channel. The present note is the $\gamma=0$ corner of
\eqref{eq:orbit}, where the orbit degenerates to a \emph{pair} and the two channels
separate completely by parity.

\section{The theorem}\label{sec:thm}

\begin{theorem}[sign structure at a real zero]\label{thm:sign}\tier{T1}
Let $f$ be real and supported in $[-L,L]$, and let a real zero pair
$\{\tfrac12+\delta,\;\tfrac12-\delta\}$, $\delta\in(0,\tfrac12]$, be present in the zero
multiset. Its contribution to the finite-window Weil form is
\[
2\left(\int f(u)e^{\delta u}du\right)\left(\int f(u)e^{-\delta u}du\right)
=\begin{cases}
+2\,\Ad(f)^2, & f \text{ even},\quad \Ad(f)=\displaystyle\int f(u)\cosh(\delta u)\,du,\\[4pt]
-2\,\Sd(f)^2, & f \text{ odd},\quad\ \Sd(f)=\displaystyle\int f(u)\sinh(\delta u)\,du,
\end{cases}
\]
for every $L>0$ and in every basis. In particular the perturbation is rank one, positive
semidefinite on the even sector and negative semidefinite on the odd sector.
\end{theorem}

\begin{proof}
Four lines; written out with every step in Appendix~\ref{app:proof}. The zero side is
$\sum_\rho h(z_\rho)$ with $h(z)=F(z)F(-z)$, $\rho=\tfrac12+iz$, and the autocorrelation
\emph{factorises} --- that is the whole content. The functional equation forces the partner,
so the pair sits at $z=\mp i\delta$ and contributes
$h(i\delta)+h(-i\delta)=2h(i\delta)=2F(i\delta)F(-i\delta)$. With
$F(i\delta)=\int f(x)e^{-\delta x}dx$ and
$e^{-\delta x}=\cosh(\delta x)-\sinh(\delta x)$, parity gives
$F(\pm i\delta)=(\Ad,\Ad)$ for even $f$ and $(-\Sd,+\Sd)$ for odd $f$. Even $f$ gives
$A\cdot A$, odd $f$ gives $S\cdot(-S)$. The sign is not a convention: it is the parity of
$F$ at a purely imaginary argument.
\end{proof}

\begin{corollary}[the even sector is anti-detecting]\label{cor:even}\tier{T1}
Suppose every non-trivial zero of $L(s,\chi)$ lies on the critical line except a single
real pair $\tfrac12\pm\delta$. Then for every $L>0$ and every real even $f$ supported in
$[-L,L]$,
\[
Q_\chi(f)=\sum_{\text{on-line }\rho}|F(\gamma_\rho)|^2+2\Ad(f)^2\ \ge\ 2\Ad(f)^2\ \ge\ 0 .
\]
Hence even-sector finite-window Weil positivity is compatible with a real zero at every
$\delta\in(0,\tfrac12)$, and no even-sector positivity computation --- at any window
length, in any basis, to any precision --- carries information about the Siegel
configuration. Equivalently, as a perturbation of the matrix: the response is $+2aa^{\mathsf T}$
with $a_k=\int\varphi_k\cosh(\delta u)\,du/\|\varphi_k\|$, so
$\lambda_{\min}(Q+2aa^{\mathsf T})\ge\lambda_{\min}(Q)$ for every basis and every $N$.
\end{corollary}

\begin{remark}\tier{T1}
This is stronger than ``the $\delta^2$ coefficient vanishes''. That weaker statement is
true --- $K(L,0)=\int_{-L}^{L}u^2\sin^2(0\cdot u)du=0$, i.e.\ $F'(0)=0$ for even $f$ ---
but it only kills the \emph{leading} coefficient. Theorem~\ref{thm:sign} is exact in
$\delta$: the response is a perfect square of the wrong sign to all orders. The even
sector is not blind at $\gamma=0$; a planted real zero makes the form \emph{more} positive.
\end{remark}

\begin{corollary}[the odd sector is maximally sensitive]\label{cor:odd}\tier{T1}
For real odd $f$ the same identity reads
$Q_\chi(f)=\sum_{\text{on-line}}|\widetilde F(\gamma_\rho)|^2-2\Sd(f)^2$, which can be
negative. Moreover the derivative-evaluation extremal in the odd sector at $\gamma=0$ is
\[
2\sup\left\{\left(\int_{-L}^{L} u\,g(u)\,du\right)^{\!2} : g \text{ odd},\ \|g\|_2=1\right\}
=2\,K_{\mathrm{odd}}(L,0)=2\cdot\frac{2L^3}{3},
\]
attained at $g\propto u$, which is the $\gamma$-uniform \emph{maximum} of the identity
$K_{\mathrm{even}}(L,\gamma)+K_{\mathrm{odd}}(L,\gamma)=2L^3/3$ together with
$K_{\mathrm{even}}(L,0)=0$.
\end{corollary}

So $\gamma=0$ is a blind spot of one \emph{sector}, not of the window, and the
complementary sector is there as sensitive as it ever gets at any ordinate.

\begin{remark}[general $f$: lowering the form needs an odd component]\label{rem:mixed}
\tier{T1} For real $f=f_{\mathrm{e}}+f_{\mathrm{o}}$ the cross integrals vanish by parity,
so the computation of Appendix~\ref{app:proof} gives
$F(i\delta)=\Ad(f_{\mathrm{e}})-\Sd(f_{\mathrm{o}})$ and
$F(-i\delta)=\Ad(f_{\mathrm{e}})+\Sd(f_{\mathrm{o}})$, and the real pair contributes
$2\bigl(\Ad(f_{\mathrm{e}})^2-\Sd(f_{\mathrm{o}})^2\bigr)$. A real pair can therefore lower
the form only through a test function with a nonzero odd component, and only when
$|\Sd(f_{\mathrm{o}})|>|\Ad(f_{\mathrm{e}})|$. The statements below are in this sense about
the odd sector itself and not about a choice of basis.
\end{remark}

\paragraph{$\gamma=0$ is the only ordinate where the sectors separate.}
The identity $K_{\mathrm{even}}+K_{\mathrm{odd}}=2L^3/3$ holds for every $\gamma$, but the
split is extreme only at $\gamma=0$. At large $\gamma$ the two are nearly equal: the
companion note measures $K_{\mathrm{even}}/K_{\mathrm{odd}}=1.00634$ at $L=2$ and
$\gamma=\gamma_1$, the first off-line ordinate of the Davenport--Heilbronn function,
$\gamma_1\approx85.7$~\cite[\S6]{DHnote}. The complete separation of roles is a
$\gamma=0$ phenomenon.

Verified at $\gamma=0$, with $K_{\mathrm{odd}}$ obtained by quadrature rather than by
subtraction:

%% generated by note/make_tables.py from the data files
%% DO NOT EDIT BY HAND
\begin{center}\small
\begin{tabular}{@{}lllll@{}}
\toprule
$L$ & $K_{\mathrm{even}}(L,0)$ & $K_{\mathrm{odd}}(L,0)$ & $2L^3/3$ & $|{\rm err}|$ \\
\midrule
$0.5$ & $0$ & $0.0833333333333$ & $0.0833333333333$ & $0$ \\
$0.8$ & $0$ & $0.341333333333$ & $0.341333333333$ & $0$ \\
$1.0$ & $0$ & $0.666666666667$ & $0.666666666667$ & $0$ \\
$1.3$ & $0$ & $1.46466666667$ & $1.46466666667$ & $0$ \\
$2.0$ & $0$ & $5.33333333333$ & $5.33333333333$ & $0$ \\
\bottomrule
\end{tabular}
\end{center}


\section{Verification}\label{sec:verif}

Three independent checks, from \src{src/verify\_siegel\_parity.py}. All three passed, and
were re-run for this note and reproduced the committed transcript byte for byte.

\paragraph{(a) The planted object was the real \emph{pair}.} Decided from the committed
Phase-5 data rather than from a reading of the code, by recomputing stored eigenvalues both
ways. The degenerate quartet is exactly twice the pair ---
$\max|{\rm quartet}(\gamma=0)-2\cdot{\rm pair}|=4.98\times10^{-60}$ --- so the two
hypotheses differ by a factor of $2$ and the data distinguishes them: all twelve probed
rows are reproduced by the \textsc{pair}, unanimously. \textbf{No factor-of-$2$ correction
is needed anywhere in this note.} \tier{T2-N}

\paragraph{(b) Rank one, correct sign, eigenvalue $=+2A^2$ / $=-2S^2$.}
$5$ characters $\times$ $L\in\{0.8,1.0,1.4\}$ $\times$ $2$ sectors $\times$
$\delta\in\{0.1,0.25,0.375,0.45,0.49\}$, $N=10$: $150$ cells, $150$ with the predicted
sign. $A$ and $S$ are computed by quadrature on the explicit normalised eigenfunction, so
this is an independent check of the closed forms, not a restatement of them. \tier{T2-N}

\begin{center}\small
\begin{tabular}{@{}ll@{}}
\toprule
quantity & worst over the $150$ cells \\
\midrule
relative error, eigenvalue vs $+2A^2$ (even) / $-2S^2$ (odd) & $8.49\times10^{-61}$ \\
rank-one residual $\max_{ij}|B_{ij}-\lambda v_iv_j|$ & $2.49\times10^{-60}$ \\
largest non-dominant eigenvalue & $1.45\times10^{-60}$ \\
the complementary integral ($S$ on even $f$, $A$ on odd $f$) & exactly $0$ \\
even-sector eigenvalues, $75$ of $75$ & $+3.14163\ \ldots\ +6.38933$ \\
odd-sector eigenvalues, $75$ of $75$ & $-0.905634\ \ldots\ -0.00643962$ \\
$K_{\mathrm{even}}+K_{\mathrm{odd}}=2L^3/3$, $K_{\mathrm{odd}}$ by quadrature & $1.56\times10^{-61}$ \\
\bottomrule
\end{tabular}
\end{center}

%% generated by note/make_tables.py from the data files
%% DO NOT EDIT BY HAND
\begin{center}\small
\begin{tabular}{@{}llllll@{}}
\toprule
$L$ & $\delta$ & $+2A_\delta^2$ (even) & $-2S_\delta^2$ (odd) & rank-one resid. & rel.\ err. \\
\midrule
$0.8$ & $0.1$ & $3.14162843133$ & $-0.00643961708992$ & $6.22\times10^{-61}$ & $8.49\times10^{-61}$ \\
$0.8$ & $0.49$ & $3.29375936732$ & $-0.159042352281$ & $3.69\times10^{-61}$ & $1.89\times10^{-61}$ \\
$1.0$ & $0.1$ & $3.93155834205$ & $-0.0125860661447$ & $6.42\times10^{-61}$ & $1.58\times10^{-61}$ \\
$1.0$ & $0.49$ & $4.23383664019$ & $-0.315813694881$ & $6.22\times10^{-61}$ & $2.46\times10^{-61}$ \\
$1.4$ & $0.1$ & $5.52110947823$ & $-0.0345998225163$ & $6.22\times10^{-61}$ & $5.64\times10^{-61}$ \\
$1.4$ & $0.49$ & $6.38932656805$ & $-0.905634273488$ & $6.22\times10^{-61}$ & $2.58\times10^{-61}$ \\
\bottomrule
\end{tabular}
\end{center}


\noindent (The extreme $\delta$ only; all $150$ cells are in
\src{data/verify\_siegel\_parity.csv}.)

\noindent\textbf{The matrices do not depend on $\chi$ at all.} Grouping the $150$ cells by
$(L,\text{sector},\delta)$ leaves exactly \emph{one} distinct eigenvalue in each of the
$30$ groups: the response to a real zero is a property of the basis and $\delta$ only. That
is why the conclusion is uniform over the family rather than five separate measurements,
and it is why the table above has no $\chi$ column.

\paragraph{(c) Pole consistency --- the strongest check in this note.}
The $\delta=\tfrac12$ real-pair perturbation was compared entrywise with this repository's
$\zeta$-side pole term $2\,\mathrm{pole\_sign}(\text{sector})P_jP_k/\|\cdot\|$, reporting
\emph{both} differences so that the data fixes the sign. \tier{T2-N}

%% generated by note/make_tables.py from the data files
%% DO NOT EDIT BY HAND
\begin{center}\small
\begin{tabular}{@{}llllll@{}}
\toprule
sector & $L$ & pole sign & $\max|B-\mathrm{Pole}|$ & $\max|B+\mathrm{Pole}|$ & vanishes \\
\midrule
\textsc{even} & $0.5$ & $+1$ & $1.07\times10^{-61}$ & $4.084$ & B $-$ Pole \\
\textsc{even} & $0.8$ & $+1$ & $3.11\times10^{-61}$ & $6.749$ & B $-$ Pole \\
\textsc{even} & $1.3$ & $+1$ & $3.6\times10^{-61}$ & $11.95$ & B $-$ Pole \\
\textsc{even\_d} & $0.5$ & $+1$ & $1.94\times10^{-62}$ & $3.281$ & B $-$ Pole \\
\textsc{even\_d} & $0.8$ & $+1$ & $3.11\times10^{-61}$ & $5.347$ & B $-$ Pole \\
\textsc{even\_d} & $1.3$ & $+1$ & $6.22\times10^{-61}$ & $9.128$ & B $-$ Pole \\
\textsc{odd} & $0.5$ & $-1$ & $1.22\times10^{-63}$ & $0.05108$ & B $-$ Pole \\
\textsc{odd} & $0.8$ & $-1$ & $9.72\times10^{-63}$ & $0.2119$ & B $-$ Pole \\
\textsc{odd} & $1.3$ & $-1$ & $1.94\times10^{-62}$ & $0.9408$ & B $-$ Pole \\
\bottomrule
\end{tabular}
\end{center}


\noindent $B=\mathrm{Pole}$ exactly in all $9$ cells, with
$\max|B-\mathrm{Pole}|\le6.22\times10^{-61}$, while the alternative $\max|B+\mathrm{Pole}|$
is $0.05108$ to $11.95$ --- wrong by $O(1)$ to $O(10)$, so the test is decisive.
\textbf{There is no normalisation mismatch: no factor of $2$ and no factor of
$\tfrac12$.} This is the strongest check because the pole term was derived independently,
for a different purpose, years-of-work earlier in this repository's own $\zeta$ assembly,
and Theorem~\ref{thm:sign} reproduces it as the $\delta=\tfrac12$ instance of $-2\Sd^2$ in
the odd sector and of $+2\Ad^2$ in the even ones.

\begin{remark}[a sign that must not be misread]
The odd sector's $\mathrm{pole\_sign}$ is $-1$, so the odd-sector pole term carries the
sign \emph{opposite to the even-sector one}, and the real-pair perturbation reproduces that
flip exactly. It is \textbf{not} opposite in sign to the pole term itself. The distinction
matters because Corollary~\ref{cor:even} rests on this sign.
\end{remark}

\paragraph{The $\chi$-analogue is identically zero, and that is the point.} $M_\chi$
reconstructs as $\mathrm{Arch}-\mathrm{LogPi}+\mathrm{LogQ}-\mathrm{Prime}$ with residual
exactly $0$: for non-principal $\chi$ there is no pole term at all. Meanwhile the real-pair
perturbation at $\delta=\tfrac12$ is \emph{not} zero for $\chi$ (largest entry
$2.673477$).
On the $\zeta$ side a pole-like object at $z=\mp i/2$ is already present in the form; on the
$\chi$ side there is nothing of the kind to cancel against a real zero.

\section{The function-field analogue}\label{sec:ff}

\tier{T1, exact over $\mathbb{Q}$} The same configuration exists in this repository's
function-field control, where RH is Weil's theorem. There a real reciprocal pair
$\{\rho,\rho^{-1}\}$ gives the rank-two Toeplitz matrix of Theorem~6 of the $\delta^2$--$L^3$
paper~\cite{deltaL3},
\[
T_R=u w^{\mathsf T}+w u^{\mathsf T},\qquad u_i=\rho^{\,i},\quad w_i=\rho^{-i},\quad 0\le i\le R .
\]
With $J$ the flip $x_i\mapsto x_{R-i}$ one has $Ju=\rho^Rw$ and $Jw=\rho^{-R}u$, so the
parity directions are the \textbf{$\rho^R$-weighted} vectors $v_\pm=u\pm\rho^Rw$, and
\[
T_R=\frac{1}{2\rho^R}\left[v_+v_+^{\mathsf T}-v_-v_-^{\mathsf T}\right],
\qquad T_Rv_\pm=\lambda_\pm v_\pm,
\qquad \lambda_\pm=(R+1)\pm\rho^{-R}\|u\|^2 .
\]
Verified exactly in \texttt{fractions.Fraction}, no floating point, for
$\rho\in\{2,\tfrac32,\tfrac54,\tfrac73,1\}\times R\in\{1,2,4,8,16\}$ --- $25$ cells, $0$
failures: the splitting identity, both eigenvector relations, $v_+$ palindromic and $v_-$
anti-palindromic under $i\mapsto R-i$, $\lambda_+>0$, and $\lambda_-\le0$ with equality iff
$\rho=1$. The negative eigenvalue, exactly:

%% generated by note/make_tables.py from the data files
%% DO NOT EDIT BY HAND
\begin{center}\small
\begin{tabular}{@{}llll@{}}
\toprule
$\rho$ & $R=1$ & $R=2$ & $R=4$ \\
\midrule
$2$ & $-1/2$ & $-9/4$ & $-261/16$ \\
$3/2$ & $-1/6$ & $-25/36$ & $-5125/1296$ \\
$5/4$ & $-1/20$ & $-81/400$ & $-168561/160000$ \\
$7/3$ & $-16/21$ & $-1600/441$ & $-6088000/194481$ \\
$1$ & $0$ & $0$ & $0$ \\
\bottomrule
\end{tabular}
\end{center}


\noindent ($R=8$ and $R=16$ are verified in the same run; their exact rationals run to
$22$ digits in the denominator and are not printed.)

\begin{remark}[a correction, carried over]\label{rem:ffcorr}
The brief that commissioned the parent extension, and the brief that commissioned this
note, both state that the eigenvalues $(R+1)\pm\|u\|\|w\|$ are attained on $u+w$ and
$u-w$, with $u+w$ palindromic. That is \textbf{false for every $\rho\ne1$}: it holds in
exactly $5$ of the $25$ cells, which are precisely the degenerate $\rho=1$ ones, and
$T_R(u+w)$ is parallel to $u+w$ only when $\rho^{2R}=1$. The missing factor is $\rho^R$.
The \emph{polarisation identity}
$T_R=\tfrac12[(u+w)(u+w)^{\mathsf T}-(u-w)(u-w)^{\mathsf T}]$ is nevertheless exact in all
$25$ cells. Neither correction touches any result of~\cite{deltaL3}, whose proof derives
the eigenvalues basis-independently and never names $u\pm w$ as eigenvectors.
\end{remark}

\paragraph{What the control shows.} The $J$-even direction carries the positive rank-one
piece and the $J$-odd direction the negative one: the same signed split, with the same
attachment of sign to parity, as $+2\Ad^2$ / $-2\Sd^2$. So the sign structure is intrinsic
to ``a real pair in a finite window'' and is not an artefact of the archimedean bookkeeping
of $L(s,\chi)$. This is a check at a level the leading-constant dictionary of~\cite{deltaL3}
did not reach: that dictionary matched a $\delta^2$ coefficient between the two sides,
whereas what matches here is the \emph{sign attached to a parity class}, an exact and
discrete datum. It does \textbf{not} transfer the consequence of
Corollary~\ref{cor:even}: in the function field RH is a theorem, so there is no
configuration to exclude.

\section{Consequences for even-sector techniques}\label{sec:lit}

\tier{T3} This section is a scope statement about what certain methods can detect. It is
not a criticism of any result, and none of the work cited claims otherwise.

Theorems~7, 8 and~11 of the $\delta^2$--$L^3$ paper~\cite{deltaL3} are labelled
even-sector in their own statements, and that paper says so globally. By
Corollary~\ref{cor:even}, the even sector is precisely the one that cannot respond to a
real zero: a positivity statement there is compatible with a real zero at any
$\delta\in(0,\tfrac12)$. So each of those three is silent on the Siegel configuration
\emph{by construction} --- not for want of window length, basis size or precision, but
because the perturbation it would have to see is positive semidefinite in that sector. The
same applies to the four-mode certified witness of Theorem~11, whose witness consists of
even quarter-wave modes.

The Connes--van Suijlekom pipeline~\cite{CvS} enters the same way at one specific point:
its hypothesis concerns a simple lowest eigenvalue with an \emph{even} eigenfunction, so
the quantity it is about is an even-sector ground state. Corollary~\ref{cor:even} says
that quantity cannot move in response to a planted real zero.

\paragraph{Zhu's certificate is a different case, and this note does not settle it.}
It would be convenient to add Zhu's $\lambda^*(0.8)\ge8.9\times10^{-18}$~\cite{Zhu} to the
list. That does \textbf{not} follow, and nothing here establishes it either way, because
the \emph{domain} of Zhu's infimum is deliberately left open by~\cite{deltaL3}, which says
only that a certificate proved on the full admissible space may be quoted unchanged. The
two cases differ in kind:
\begin{itemize}\itemsep2pt
\item if $\lambda^*(L)$ is an infimum over the \textbf{even sector only}, then by
Corollary~\ref{cor:even} it carries no information about a real zero, exactly as for
Theorems~7, 8 and~11;
\item if it is an infimum over the \textbf{full admissible space}, then it \emph{does}
constrain one, since positivity on the full space includes the odd sector where the
response is $-2\Sd(f)^2$ --- and that would be the opposite conclusion, and the more
interesting one.
\end{itemize}
Which case holds is a property of a paper that was not obtained or read for this work. \textbf{No
claim is made in either direction.} The brief that commissioned this note listed Zhu's
certificate among the even-sector cases; that listing is not supported, and is corrected
here rather than repeated.

\paragraph{Where a window-based statement about real zeros would have to come from.}
From a test function with a nonzero odd component, and from nothing else: by
Remark~\ref{rem:mixed} a real pair contributes
$2(\Ad(f_{\mathrm{e}})^2-\Sd(f_{\mathrm{o}})^2)$, so lowering the form requires the odd part
to dominate. The odd sector at small $L$ is where that is easiest, the response there being
negative and largest relative to the ordinate (Corollary~\ref{cor:odd}). But the route is subject to the obstruction the companion
note~\cite{DHnote} measured on the Davenport--Heilbronn function: there, the negativity
produced by an off-line orbit was partly cancelled by the remaining spectrum, and turning a
detector response into an exclusion needs an unconditional lower bound on the
remaining-zero form. No such bound is known to the author for any Dirichlet window, and
none is produced here (non-claim~4) --- a statement about what is available, not a proof
that none exists.

\section{Illustration}\label{sec:illus}

\tier{T2-N; every row} Smallest scanned $L$ at which a planted real zero at each $\delta$
produces a negative odd-sector Rayleigh quotient, at fixed basis dimension $N=12$, over
$L\in\{0.5,0.8,1.1,1.4,1.7,2.0,2.3,2.6,3.0\}$ ($270$ rows, $228$ detections):

%% generated by note/make_tables.py from the data files
%% DO NOT EDIT BY HAND
\begin{center}\small
\begin{tabular}{@{}llllll@{}}
\toprule
 & $\delta=0.1$ & $\delta=0.25$ & $\delta=0.375$ & $\delta=0.45$ & $\delta=0.49$ \\
\midrule
$\zeta$ & $0.5$ & $0.5$ & $0.5$ & $0.5$ & $0.5$ \\
$\chi_{-3}$ & $0.8$ & $0.8$ & $0.8$ & $0.8$ & $0.8$ \\
$\chi_{-4}$ & $1.1$ & $1.1$ & $1.1$ & $0.8$ & $0.8$ \\
$\chi_{5}$ & $1.1$ & $1.1$ & $0.8$ & $0.8$ & $0.8$ \\
$\chi_{-7}$ & $1.4$ & $1.1$ & $1.1$ & $1.1$ & $1.1$ \\
$\chi_{8}$ & $1.4$ & $1.1$ & $1.1$ & $1.1$ & $1.1$ \\
\bottomrule
\end{tabular}
\end{center}


\noindent So: \textbf{odd sector, by $L=1.4$ with a $12$-dimensional basis, for every
character and every $\delta\ge0.1$}; as early as $L=0.5$ for $\zeta$ and $L=0.8$ for
$\chi_{-3}$. And \textbf{even sector: never}, at any $L$, any $N$, any $\delta$ --- in the
$180$ even-sector rows of the Phase-5 scan, $180$ have response $\ge0$ and $0$ are
detections, while all $180$ odd-sector rows have response $<0$ and $135$ are detections.
That is Theorem~\ref{thm:sign} as a prediction, confirmed, not as evidence for itself.

\begin{remark}[two caveats that change how the table should be read]
First, the $\zeta$ row is nearly vacuous: the odd-sector baseline there is
$\lambda^*_{\mathrm{odd}}(\zeta,0.8)=4.98077\times10^{-14}$, so $\zeta$'s odd form is
already within $10^{-14}$ of singular and \emph{any} negative rank-one perturbation
detects. The informative rows are $\chi_{-4}$ and $\chi_5$ at $L=0.8$, which clear
baselines of $0.0988$ and $0.0679$ respectively (both at $N=12$). Second, the ordering of the onsets is the
ordering of the odd-sector baselines, and across these five characters the conductor $q$ and
the first ordinate $\gamma_1$ are not independently varied, so the ordering is not evidence
of a conductor effect.
\end{remark}

\section{What this rules out, and what it does not}

\paragraph{Ruled out.} Reading any even-sector finite-window positivity result --- at any
$L$, in any basis, to any precision --- as evidence against a real zero near $1$. By
Theorem~\ref{thm:sign} the perturbation such a result would have to detect is positive
semidefinite in that sector, so the computation is not insensitive to the configuration but
\emph{anti-}sensitive to it: planting the zero improves the very quantity being certified.

\paragraph{Not ruled out, and not claimed.} Whether any real zero near $1$ exists, for
$\zeta$ or for any $L(s,\chi)$: nothing here bears on it. An odd-sector certificate: none is
given, and none is known for any Dirichlet window, so the odd-sector results are
detector-sensitivity measurements and not exclusions. Complex characters: the whole analysis
assumes $\chi$ real, since it is the reality of $\chi$ that makes $M_\chi$ real symmetric
and the parity split available. And the masking obstruction of the companion note is
untouched --- this note says where a real zero would be visible, not that it could be
isolated from the rest of the spectrum once it was.

\appendix

\section{The proof, with the factorisation written out}\label{app:proof}

The content of Theorem~\ref{thm:sign} is that the autocorrelation factorises. Writing it
out once, with every step:
\begin{align*}
\text{zero side} &= \sum_\rho h(z_\rho), \qquad h(z)=F(z)F(-z),\qquad \rho=\tfrac12+iz,\\
\text{real pair } \{\tfrac12\pm\delta\} &\ \longrightarrow\ z=\mp i\delta
  &&\text{(the functional equation forces the partner)}\\
\text{contribution} &= h(i\delta)+h(-i\delta)=2h(i\delta)
  &&\text{($h$ even)}\\
 &= 2\,F(i\delta)\,F(-i\delta)
  &&\text{(the factorisation --- the whole content)}\\
F(i\delta) &= \int f(x)e^{-\delta x}dx
  = \int f\cosh(\delta x)dx-\int f\sinh(\delta x)dx\\
f \text{ even} &:\quad \Ad-0=\Ad,\qquad F(-i\delta)=\Ad+0=+\Ad\\
f \text{ odd\ } &:\quad 0-\Sd=-\Sd,\qquad F(-i\delta)=0+\Sd=+\Sd
\end{align*}
so the product is $+\Ad^2$ for even $f$ and $-\Sd^2$ for odd $f$, and the contribution is
twice it. In matrix form, with $a_k=\int\varphi_k\cosh(\delta u)du/\|\varphi_k\|$ and
$s_k=\int\psi_k\sinh(\delta u)du/\|\psi_k\|$ on the even and odd bases, the perturbations
are $+2aa^{\mathsf T}$ and $-2ss^{\mathsf T}$: rank one, of the stated signs.

Two points worth isolating. The sign is \emph{not} a convention: it is the parity of $F$ at
a purely imaginary argument, and $F(\pm i\delta)$ are both real for real $f$ of definite
parity, with the \emph{same} sign in the even case and \emph{opposite} signs in the odd
case. And the statement is exact in $\delta$: nothing is expanded, so there is no
$\delta\to0$ regime in which it fails.

\section{Retraction log}\label{app:retr}

Reproduced from the parent extension's \S0.2, verbatim in substance, plus this session's.

\begin{description}\itemsep3pt
\item[R1.] \src{src/regression\_zeta.py} asserted the $\zeta$-side values against the
parent's committed CSVs at relative tolerance $10^{-20}$ and reported $14$ failures.
\textbf{Retracted: the tolerance was wrong, the values are right.} The parent prints
eigenvalues to twelve significant digits, so a full-precision recomputation can agree with
the stored \emph{decimal string} only to about $10^{-12}$; every ``failure'' sat at
$4\times10^{-15}\ldots1.7\times10^{-12}$. Tolerance corrected to $5\times10^{-12}$. No
computed value changed. The binding precision on any regression against this repository's
CSVs is the CSVs' own.

\item[R2.] \src{src/phase2\_gateA\_primeflip.py} built $L=\texttt{mpf('0.8')}$ \emph{before}
setting the working precision, so $L$ carried only $15$ digits into a $40$-digit run.
\textbf{Fixed.} The affected quantities were control ratios of order $10^4$, so no reported
conclusion moved, but the ordering was wrong and is recorded rather than quietly repaired.

\item[R3.] \src{src/planting.py::quartet\_block} initially omitted the factor
$h_{\rm sign}=-1$ for the odd basis, where the repository's $F$ returns $F/i$.
\textbf{Fixed before use}, caught by the planting self-test ($\mathrm{quartet}(\gamma,0)=
2\cdot\mathrm{pair}(\gamma)$ failed by $4.04$). No result was computed with the bad
version.

\item[R4.] A commit message listed a fourth retraction --- the $\zeta$ Connes window ---
but the row was never added to the table. \textbf{Fixed.} The substance: that file's fourth
window is $L=\log(13)/2$ stored rounded to $10$ digits, and $\lambda_{\min}\approx10^{-29}$
there is hypersensitive because $2L$ sits exactly at the $n=13$ prime-power cutoff, so
reconstructing $L$ from the stored string moved the answer by $1.8\times10^{-8}$ relative.
\emph{The documentation gap is itself the retraction: a commit message is not the record.}

\item[R5.] \src{src/verify\_siegel\_parity.py} first asserted a stored column at relative
tolerance $10^{-25}$, which would have reported $12$ spurious failures at observed
deviations of $4\times10^{-11}\ldots1.3\times10^{-10}$. \textbf{Fixed before the run that
produced the committed transcript}: the stored column carries ten significant digits, so
the binding tolerance is the file's. \emph{This is R1 recurring --- the lesson did not
transfer the first time.}

\item[R6 (this session).] Two claims in the brief that commissioned this note were repeated
from an earlier brief after the parent extension had already refuted them, and are
\textbf{not} asserted here: that the function-field eigenvalues are attained on $u\pm w$
with $u+w$ palindromic (false for every $\rho\ne1$; Remark~\ref{rem:ffcorr}), and that
Zhu's certificate is one of the even-sector cases that carry no information about the
Siegel configuration (not established either way; \S\ref{sec:lit}). Both are corrected in
place rather than carried forward.
\end{description}

\section{Reproduction}\label{app:repro}

Code and data: \url{https://github.com/innerlightr-wq/weil-window-control}, directory
\src{extensions/dirichlet-windows-2026-10-23/}. Python~3.14.7 (CPython),
\texttt{mpmath}~1.3.0; no \texttt{numpy}, no \texttt{float64} on any path to a reported
number; \S\ref{sec:ff} uses \texttt{fractions.Fraction} only.

\begin{center}\small
\begin{tabular}{@{}lll@{}}
\toprule
Result & Script & Data \\
\midrule
\S\ref{sec:verif} checks (a), (b), (c) & \src{src/verify\_siegel\_parity.py} & \src{data/verify\_siegel\_parity.csv}, \src{results/verify\_siegel\_parity.json} \\
\S\ref{sec:thm} the $K$ identity & \src{src/phase4b\_sector\_extremal.py} & \src{data/phase4b\_sector\_extremal.csv} \\
\S\ref{sec:ff} function field & \src{src/ff\_parity\_split.py} & \src{results/ff\_parity\_split.json} \\
\S\ref{sec:illus} onset, odd sector & \src{src/phase5\_onset.py} & \src{data/phase5\_onset\_odd.csv} \\
\S\ref{sec:illus} even/odd scan & \src{src/phase5\_siegel.py} & \src{data/phase5\_siegel.csv} \\
\bottomrule
\end{tabular}
\end{center}

\noindent \texttt{make verify}, in the note directory, runs
\src{note/verify\_note\_numbers.py}: it rebuilds every quantity printed above from the data
files, scans this document for numeric literals, fails if any literal is not derived or is
printed to more significant digits than its source carries, regenerates the five
machine-generated tables and compares them, and re-checks that the three prerequisite
verifications are recorded as passed. \texttt{make} builds the PDF. The full extension
suite is \src{./run\_suite.sh}.

\section*{Licence}
\addcontentsline{toc}{section}{Licence}

As the parent repository: code MIT, this note and the data CC~BY~4.0. CC~BY~4.0 requires
attribution; if you use the note or the data, please cite it by its DOI and the parent
paper~\cite{deltaL3} separately.

\section*{Acknowledgments}
\addcontentsline{toc}{section}{Acknowledgments}

AI assistance was used for algebra checking, numerical verification, adversarial review,
literature organization, and manuscript drafting. All conceptual development, research
direction, interpretation, and final responsibility for the content belong to the author.

\begin{thebibliography}{9}

\bibitem{CvS} A. Connes and W. D. van Suijlekom,
\emph{Quadratic forms, real zeros and echoes of the spectral action}, arXiv:2511.23257.

\bibitem{DHnote} E. De Jes\'us,
\emph{A Certified Negative Finite-Window Weil Witness for the Davenport--Heilbronn
Function, with the Remaining-Zero Form Positive}, Zenodo, 2026.
\href{https://doi.org/10.5281/zenodo.23608288}{doi:10.5281/zenodo.23608288}.

\bibitem{deltaL3} E. De Jes\'us,
\emph{A $\delta^2$--$L^3$ Sensitivity Law for Finite-Window Weil Positivity, with a
Function-Field Control of the Connes--van Suijlekom Pipeline}, Zenodo, 2026.
\href{https://doi.org/10.5281/zenodo.23250932}{doi:10.5281/zenodo.23250932}.

\bibitem{HP19} E. Hallouin and M. Perret,
\emph{A unified viewpoint for upper bounds for the number of points of curves over finite
fields via Euclidean geometry and semi-definite symmetric Toeplitz matrices},
Trans. Amer. Math. Soc. \textbf{372} (2019), no.~8, 5409--5451.

\bibitem{Zhu} X. Zhu,
\emph{Weil positivity in compact windows: a finite reduction, certified two-sided bounds,
and a Landau--Widom decay law}, arXiv:2608.24827.

\end{thebibliography}

\end{document}
